\begin{appendix}
\onecolumn
\nolinenumbers
\section{Funcionamiento del código, bloque por bloque}
\label{ap:codigo}

Este apéndice recorre \texttt{src/poisson\_comovil\_cpu.c} completo, dividido en nueve bloques. Son los mismos bloques en que el cuaderno \texttt{poisson\_comovil\_LFIS428.ipynb} escribe el archivo, una celda por bloque, de modo que el código del cuaderno y el del repositorio son idénticos. Para cada bloque se explica qué hace, por qué, y cómo se trasladará a CUDA.

\subsection{Bloque 1: cabecera, bibliotecas y tipos}

El comentario inicial documenta la física (la ecuación de Poisson comóvil en unidades de código, $\nabla^2\tilde\phi=\tfrac32(\Omega_m/a)\,\delta$), los cambios respecto al código de referencia y cómo compilar y ejecutar. Se incluyen \texttt{math.h} (funciones trigonométricas y \texttt{sqrt}), \texttt{stdio.h} (impresión y archivos), \texttt{stdlib.h} (memoria dinámica y \texttt{atoi}/\texttt{atof}) y \texttt{sys/time.h} (reloj de pared). La línea \texttt{typedef double real} define un único tipo de punto flotante para todo el programa: cambiarlo a \texttt{float} es una de las optimizaciones posibles en GPU, a costa de limitar la tolerancia alcanzable a $\sim10^{-7}$.

\begin{lstlisting}
/*
 * Proyecto: ecuación de Poisson comóvil en 2D (versión CPU de referencia)
 * Modificación de poisson_cpu.c para el potencial gravitatorio peculiar.
 *
 * Ecuación física (coordenadas comóviles x, potencial peculiar phi):
 *     lap_x phi = (3/2) Omega_m H0^2 delta / a
 *
 * Unidades de código: x~ = x / L  (caja periódica [0,1) x [0,1)),
 *                     phi~ = phi / (H0 L)^2
 *     =>  lap phi~ = (3/2) (Omega_m / a) delta        (sin constantes)
 *
 * Cambios respecto al código del profesor:
 *   1. Bordes periódicos: N puntos por lado, h = 1/N, índices con módulo.
 *      Se actualizan TODOS los puntos (no hay borde fijo).
 *   2. Fuente f = (3/2)(Om/a)(delta - <delta>): se resta el promedio para
 *      cumplir la condición de compatibilidad  sum(f) = 0.
 *   3. El potencial queda definido salvo una constante: se fija <phi> = 0.
 *   4. SOR con w_opt periódico: el modo de error más lento es (1,0),
 *      k = (2 pi, 0), con factor de Jacobi rho_J = [1 + cos(2 pi h)]/2
 *      = cos^2(pi h). Entonces w = 2 / (1 + sqrt(1 - rho_J^2))
 *      = 2 / (1 + sin(pi h) sqrt(1 + cos^2(pi h))).
 *   5. Dos fuentes posibles:
 *        modo 0: test con un solo modo, delta = A cos(2 pi m x) cos(2 pi n y)
 *                (solución discreta y continua exactas conocidas)
 *        modo 1: delta leído de 'delta.bin' (p.ej. el campo gaussiano
 *                que genera el cuaderno poisson_comovil_LFIS428.ipynb)
 *
 * Compilar:  gcc -O2 poisson_comovil_cpu.c -lm -o poisson_comovil_cpu.x
 * Ejecutar:  ./poisson_comovil_cpu.x [N] [metodo] [tol] [modo] [Om] [a]
 *            N debe ser par (rojo-negro periódico).
 * Salida:    salida_cpu.bin con dos campos: phi y delta (formato del profesor)
 */

#include <math.h>
#include <stdio.h>
#include <stdlib.h>
#include <sys/time.h>

#define PI 3.14159265358979323846

typedef double real;
\end{lstlisting}

\subsection{Bloque 2: tiempo y entrada/salida}

Tres funciones auxiliares que no tocan la física. \texttt{segundos} lee el reloj del sistema con resolución de microsegundos; se usa solo alrededor del bucle iterativo, de modo que el tiempo medido no incluye la lectura de datos ni la escritura de resultados. \texttt{escribir\_salida} guarda una cabecera de cuatro enteros $(N_x, N_y, n_{\rm campos}, \text{bytes por número})$ seguida de los campos en orden \texttt{[y][x]}; es el mismo formato del código de referencia, así que las herramientas del curso lo leen sin cambios. \texttt{leer\_delta} hace lo inverso para la fuente $\delta$ y devuelve un código de error distinto si el archivo no existe, si su tamaño no coincide con $N$ o si está incompleto.

\begin{lstlisting}
double segundos(void) {
  struct timeval tp;
  gettimeofday(&tp, NULL);
  return ((double)tp.tv_sec + (double)tp.tv_usec * 1.e-6);
}

void escribir_salida(const char *nombre, int nx, int ny, int ncampos,
                     real **campos) {
  FILE *f = fopen(nombre, "wb");
  int cabecera[4] = {nx, ny, ncampos, (int)sizeof(real)};
  fwrite(cabecera, sizeof(int), 4, f);
  for (int c = 0; c < ncampos; c++)
    fwrite(campos[c], sizeof(real), (size_t)nx * ny, f);
  fclose(f);
}

// Lee delta desde un archivo con el mismo formato (cabecera + 1 campo)
int leer_delta(const char *nombre, real *delta, int N) {
  FILE *f = fopen(nombre, "rb");
  if (!f) return 1;
  int cab[4];
  if (fread(cab, sizeof(int), 4, f) != 4 || cab[0] != N || cab[1] != N ||
      cab[3] != (int)sizeof(real)) {
    fclose(f);
    return 2;
  }
  size_t n = fread(delta, sizeof(real), (size_t)N * N, f);
  fclose(f);
  return (n == (size_t)N * N) ? 0 : 3;
}
\end{lstlisting}

\subsection{Bloque 3: vecinos periódicos}

Es el cambio central respecto al código de referencia: en la caja periódica el vecino derecho de $i=N-1$ es $i=0$ y el izquierdo de $i=0$ es $i=N-1$. Se usan comparaciones en lugar de \texttt{(i+1)\%N}, que en C da un resultado negativo para $i-1<0$ y además es una operación costosa. Los arreglos son planos: el punto $(i,j)$ está en \texttt{j*N + i}, con $i$ el índice rápido (dirección $x$). \texttt{vecinos} devuelve $u_{i+1,j}+u_{i-1,j}+u_{i,j+1}+u_{i,j-1}$, la suma que aparece tanto en la regla de actualización como en el laplaciano discreto. Al ser \texttt{static inline}, el compilador la inserta en cada bucle sin costo de llamada. En CUDA estas tres funciones pasan a ser funciones \texttt{\_\_device\_\_}.

\begin{lstlisting}
// Vecinos periódicos
static inline int mas(int i, int N) { return (i + 1 == N) ? 0 : i + 1; }
static inline int menos(int i, int N) { return (i == 0) ? N - 1 : i - 1; }

// Suma de los 4 vecinos del punto (i,j) con bordes periódicos
static inline real vecinos(const real *u, int i, int j, int N) {
  return u[j * N + mas(i, N)] + u[j * N + menos(i, N)] +
         u[mas(j, N) * N + i] + u[menos(j, N) * N + i];
}
\end{lstlisting}

\subsection{Bloque 4: los métodos de relajación}

\texttt{iteracion\_jacobi} aplica la regla $u_{i,j}\leftarrow\tfrac14\left(\sum_{\rm vecinos}u - h^2 f_{i,j}\right)$ a todos los puntos, leyendo de \texttt{u} y escribiendo en \texttt{u\_new}: como ningún punto lee un valor ya actualizado, cada punto es independiente y la función se traduce directamente en un kernel con un \emph{thread} por punto. \texttt{media\_iteracion\_rb} actualiza solo los puntos con $(i+j)\bmod 2 = \texttt{color}$, en el mismo arreglo. El primer $i$ de cada fila es \texttt{(j+color)\%2}, y el paso de 2 recorre el resto del color. La actualización $u\leftarrow u+\omega\,(u^{\rm GS}-u)$ es Gauss-Seidel con $\omega=1$ y SOR con $1<\omega<2$. Los puntos de un color solo dependen de los del otro, de modo que dentro de cada media iteración también son independientes; eso exige $N$ par en la caja periódica.

\begin{lstlisting}
void iteracion_jacobi(const real *u, real *u_new, const real *f, int N,
                      real h2) {
  for (int j = 0; j < N; j++)
    for (int i = 0; i < N; i++)
      u_new[j * N + i] = 0.25 * (vecinos(u, i, j, N) - h2 * f[j * N + i]);
}

// Media iteración rojo-negro sobre los puntos con (i + j) % 2 == color
void media_iteracion_rb(real *u, const real *f, int N, real h2, real w,
                        int color) {
  for (int j = 0; j < N; j++)
    for (int i = (j + color) % 2; i < N; i += 2) {
      int idx = j * N + i;
      real gs = 0.25 * (vecinos(u, i, j, N) - h2 * f[idx]);
      u[idx] += w * (gs - u[idx]);
    }
}
\end{lstlisting}

\subsection{Bloque 5: reducciones sobre la malla}

Cuatro operaciones que combinan toda la malla en un número (reducciones). \texttt{norma\_residuo} calcula $\|r\|$ con $r=f-\nabla_h^2u$ en cada punto: es el criterio de parada. \texttt{norma} da $\|f\|$, que normaliza el residuo. \texttt{promedio} y \texttt{restar\_promedio} implementan la física de la caja periódica: la condición de compatibilidad $\sum f=0$ y la fijación del modo cero $\langle u\rangle=0$. En GPU las reducciones no son triviales, porque muchos \emph{threads} suman sobre la misma variable; se implementan con una reducción en árbol en memoria compartida o con bibliotecas como CUB o Thrust.

\begin{lstlisting}
real norma_residuo(const real *u, const real *f, int N, real h2) {
  real suma = 0.0;
  for (int j = 0; j < N; j++)
    for (int i = 0; i < N; i++) {
      real lap = (vecinos(u, i, j, N) - 4.0 * u[j * N + i]) / h2;
      real r = f[j * N + i] - lap;
      suma += r * r;
    }
  return sqrt(suma);
}

real norma(const real *f, int N) {
  real suma = 0.0;
  for (int k = 0; k < N * N; k++) suma += f[k] * f[k];
  return sqrt(suma);
}

real promedio(const real *f, int N) {
  real suma = 0.0;
  for (int k = 0; k < N * N; k++) suma += f[k];
  return suma / ((real)N * N);
}

// Fija el modo cero: <u> = 0
void restar_promedio(real *u, int N) {
  real m = promedio(u, N);
  for (int k = 0; k < N * N; k++) u[k] -= m;
}
\end{lstlisting}

\subsection{Bloque 6: parámetros, discretización y memoria}

\texttt{main} lee hasta seis argumentos con valores por defecto: $N$, el método (0 Jacobi, 1 Gauss-Seidel rojo-negro, 2 SOR rojo-negro), la tolerancia, el modo de la fuente, $\Omega_m$ y $a$. Valida que el método exista y que $N$ sea par. Calcula $h=1/N$ (en la caja periódica hay $N$ puntos distintos, $x_i=ih$) y el $\omega$ óptimo periódico $\omega=2/\big(1+\sin\pi h\sqrt{1+\cos^2\pi h}\big)$, que viene del modo de error más lento, $(1,0)$. \texttt{chequeo} fija cada cuántas iteraciones se calcula el residuo. Reserva los cuatro campos de $N^2$ números: \texttt{u} y \texttt{u\_new} con \texttt{calloc}, de modo que la iteración parte de $u=0$, y \texttt{f} y \texttt{delta} con \texttt{malloc}.

\begin{lstlisting}
int main(int argc, char **argv) {
  int N = (argc > 1) ? atoi(argv[1]) : 256;
  int metodo = (argc > 2) ? atoi(argv[2]) : 2;
  real tol = (argc > 3) ? atof(argv[3]) : 1e-8;
  int modo = (argc > 4) ? atoi(argv[4]) : 0;
  real Om = (argc > 5) ? atof(argv[5]) : 0.3;
  real a = (argc > 6) ? atof(argv[6]) : 1.0;

  const char *nombres[3] = {"Jacobi", "Gauss-Seidel rojo-negro",
                            "SOR rojo-negro"};
  if (metodo < 0 || metodo > 2 || N % 2 != 0) {
    fprintf(stderr, "metodo debe ser 0, 1 o 2 y N debe ser par\n");
    return 1;
  }

  real h = 1.0 / N;  // periódico: x_i = i h, i = 0..N-1
  real h2 = h * h;
  // rho_J = cos^2(pi h) (modo (1,0)) => w = 2/(1+sqrt(1-rho_J^2))
  real w = (metodo == 2)
               ? 2.0 / (1.0 + sin(PI * h) * sqrt(1.0 + pow(cos(PI * h), 2)))
               : 1.0;
  int chequeo = 100;
  int max_iter = 10000000;

  real *u = (real *)calloc((size_t)N * N, sizeof(real));
  real *u_new = (real *)calloc((size_t)N * N, sizeof(real));
  real *f = (real *)malloc((size_t)N * N * sizeof(real));
  real *delta = (real *)malloc((size_t)N * N * sizeof(real));
\end{lstlisting}

\subsection{Bloque 7: la fuente comóvil}

Construye $\delta$ de dos maneras. En el modo 0 es un único modo de Fourier, $\delta=\cos(2\pi mx)\cos(2\pi ny)$ con $(m,n)=(1,2)$, cuya solución discreta y continua se conocen exactamente. En el modo 1 se lee desde \texttt{delta.bin}, por ejemplo un campo gaussiano aleatorio. Luego forma la fuente en unidades de código, $f=\tfrac32(\Omega_m/a)\,(\delta-\langle\delta\rangle)$: restar el promedio impone $\sum f=0$, sin lo cual la ecuación periódica no tiene solución. Calcula $\|f\|$ para normalizar el residuo e imprime los parámetros de la corrida.

\begin{lstlisting}
  // Parámetros del test de un modo
  const int m = 1, n = 2;
  const real A = 1.0;

  if (modo == 0) {
    for (int j = 0; j < N; j++)
      for (int i = 0; i < N; i++)
        delta[j * N + i] =
            A * cos(2 * PI * m * i * h) * cos(2 * PI * n * j * h);
  } else {
    int e = leer_delta("delta.bin", delta, N);
    if (e) {
      fprintf(stderr, "error %d leyendo delta.bin (¿mismo N?)\n", e);
      return 1;
    }
  }

  // Fuente comóvil en unidades de código, con promedio restado
  real dmed = promedio(delta, N);
  real pref = 1.5 * Om / a;
  for (int k = 0; k < N * N; k++) f[k] = pref * (delta[k] - dmed);
  real norma_f = norma(f, N);

  printf("Poisson comóvil 2D periódico: N = %d, método = %s", N,
         nombres[metodo]);
  if (metodo == 2) printf(" (w = %.6f)", w);
  printf(", tol = %.1e, Om = %.3f, a = %.3f, <delta> = %.2e\n", tol, Om, a,
         dmed);
\end{lstlisting}

\subsection{Bloque 8: el bucle iterativo}

El corazón del programa. Repite bloques de \texttt{chequeo}$=100$ iteraciones: en Jacobi, una iteración y un intercambio de punteros entre \texttt{u} y \texttt{u\_new} (no se copian datos); en Gauss-Seidel y SOR, la media iteración roja seguida de la negra. Después de cada bloque resta el promedio de $u$, porque el modo cero no está fijado por la ecuación y el redondeo lo haría derivar, y calcula el residuo relativo $\|f-\nabla_h^2u\|/\|f\|$. Termina cuando cae bajo la tolerancia. El reloj encierra solo este bucle. En la versión CUDA todo este bloque corre en la GPU y solo el número del residuo vuelve al host, una vez por bloque.

\begin{lstlisting}
  double t0 = segundos();
  int iter = 0;
  real res = 1.0;
  while (iter < max_iter) {
    for (int k = 0; k < chequeo; k++) {
      if (metodo == 0) {
        iteracion_jacobi(u, u_new, f, N, h2);
        real *tmp = u;
        u = u_new;
        u_new = tmp;
      } else {
        media_iteracion_rb(u, f, N, h2, w, 0);
        media_iteracion_rb(u, f, N, h2, w, 1);
      }
    }
    iter += chequeo;
    restar_promedio(u, N);  // evita que el modo cero derive por redondeo
    res = norma_residuo(u, f, N, h2) / norma_f;
    if (res < tol) break;
  }
  double t_total = segundos() - t0;

  printf("Iteraciones: %d, residuo relativo = %.3e\n", iter, res);
  printf("Tiempo de cómputo: %.3f s (%.3f ms por iteración)\n", t_total,
         1e3 * t_total / iter);
\end{lstlisting}

\subsection{Bloque 9: validación, salida y limpieza}

En el modo 0 compara la solución con dos referencias exactas. La discreta usa el autovalor del laplaciano de cinco puntos, $\lambda_d=-\tfrac{4}{h^2}\left[\sin^2(\pi mh)+\sin^2(\pi nh)\right]$, y mide el error del \emph{solver}; la continua usa $\lambda_c=-4\pi^2(m^2+n^2)$ y mide el error de discretización, que debe escalar como $h^2$. Ambos errores se normalizan por el máximo de la solución. Finalmente escribe $\tilde\phi$ y $\delta$ en \texttt{salida\_cpu.bin} y libera la memoria.

\begin{lstlisting}
  if (modo == 0) {
    // Autovalores del laplaciano discreto y continuo de cos(2pi m x)cos(2pi n y)
    real lam_d =
        -4.0 / h2 * (pow(sin(PI * m * h), 2) + pow(sin(PI * n * h), 2));
    real lam_c = -4.0 * PI * PI * (m * m + n * n);
    real err_d = 0.0, err_c = 0.0, umax = 0.0;
    for (int j = 0; j < N; j++)
      for (int i = 0; i < N; i++) {
        real modo_ij = cos(2 * PI * m * i * h) * cos(2 * PI * n * j * h);
        real ud = pref * A * modo_ij / lam_d;
        real uc = pref * A * modo_ij / lam_c;
        real v = u[j * N + i];
        if (fabs(v - ud) > err_d) err_d = fabs(v - ud);
        if (fabs(v - uc) > err_c) err_c = fabs(v - uc);
        if (fabs(uc) > umax) umax = fabs(uc);
      }
    printf("Error vs. solución discreta exacta (solver):        %.3e\n",
           err_d / umax);
    printf("Error vs. solución continua (discretización, ~h^2): %.3e "
           "(h^2 = %.3e)\n",
           err_c / umax, h2);
  }

  real *campos[2] = {u, delta};
  escribir_salida("salida_cpu.bin", N, N, 2, campos);

  free(u);
  free(u_new);
  free(f);
  free(delta);
  return 0;
}
\end{lstlisting}

\end{appendix}
